\chapter{Navier--Stokes: what the average keeps}

\textbf{Purpose:} Hold every claim of the second observation of the thought experiment beside what
stands behind it, kept the same way as the first. \textbf{Scope:} the second observation, and the
research paper chapter \emph{What the average keeps}, whose propositions this chapter names by
number.

\section{The second observation}

Take a quintillion molecules of a fluid. Each one feels the field and every neighbor. Reducing
what they do to a few numbers per point looks absurd once that is in view. The observation does not
say the average is wrong. It says the average over enough volume and enough time gives an entropic
equivalence that is approximate and not exact, and that the equations are built on that average.

\section{Entries}

\begin{longtable}{@{}>{\raggedright\arraybackslash}p{0.40\textwidth}%
>{\raggedright\arraybackslash}p{0.12\textwidth}%
>{\raggedright\arraybackslash}p{0.40\textwidth}@{}}
\toprule
Claim & Status & What stands behind it \\
\midrule
\endhead
\bottomrule
\endfoot
The field is continuous and the molecules sit in it; the differences of potential between any two
are uncountably many.
& proved
& Proposition 8: a continuous field takes every value between its values at two molecules, and an
interval of real numbers is uncountable. \\[0.6em]
Normalized to 0 and 1, every one of those values is kept.
& proved
& Proposition 8: $t \mapsto t/(1+t)$ carries $[0,\infty)$ one to one onto $[0,1)$, keeping order,
and cuts none off. \\[0.6em]
The continuum of the first observation and the molecules of the second pull in opposite
directions, and the research paper must take one.
& not so
& Put forward while this observation was entered, and taken back. The field is defined at every
point and the molecules are a finite number of points in it. Both hold at once. Proposition 8.
\\[0.6em]
Equation (1), read as a law of a real fluid, is a law of averages over many molecules.
& proved
& The research paper reads it so in its second chapter. Equation (1) carries no molecules and
holds a few numbers per point. \\[0.6em]
The average throws away what each molecule is doing.
& proved
& Proposition 9: a cell of $N$ molecules keeps four numbers, and every state in a set of dimension
$6N-4$ gives the same four. \\[0.6em]
The average is an entropic equivalence.
& proved
& Proposition 9: the set of states with the same cell values is the equivalence class, and its size
is what the Boltzmann entropy measures. \\[0.6em]
Over enough volume the equivalence is approximate and never exact.
& proved, under independence
& Proposition 10: the scatter of the average is $s/\sqrt N$, positive for every finite $N$. A liquid
is not independent, and the rate is then the best case. \\[0.6em]
Over enough time the same holds.
& theory
& The time form, with the time over which the motion stays correlated in the place of one
molecule, is stated and not proved. \\[0.6em]
A solution that breaks down leaves too little volume and too little time to average over.
& proved
& Proposition 11: in the class of the forced construction, the length $\ell$ of the flow and the
time $\ell^2/\nu$ have no positive lower bound. \\[0.6em]
In water, that happens at lengths a solution passes through.
& computed
& 33.4 molecules in a cube of side one nanometer at $20\,^\circ$C, with a scatter of 17 percent
under independence, and $\ell^2/\nu$ of one picosecond there. Density from IAPWS R6-95 (2018)~\cite{src:IAPWS-R6-95-2018},
viscosity from IAPWS R12-08~\cite{src:IAPWS-R12-08}, both evaluated by programs that reproduce their
verification tables. \\
\end{longtable}

The status \emph{computed} is added here: a number worked out from stated inputs, with the inputs
named and their source given.

\section{What it wants}

The time form of Proposition 10 wants a proof under a stated rate at which the motion stops being
correlated with itself.

The entries on water want the collision and correlation times of liquid water, read from a source,
set against the times of Proposition 11. Where those cross is where the average stops being
available, and that number has not been taken.

The count of Proposition 9 wants extending to the field between the molecules: what an average over
a cell keeps of a continuous field, and what it drops.
